% Notes-only C.51 candidate. Canonical sources are frozen for checkpoint 17.
\subsection{An explicit finite outer detector for every multiplier}

\begin{statement}{C.51}{Open}{An explicit full-group outer detector using only the primes two and three}
Let $D$ be the explicitly specified fully characteristic finite quotient
of C.50. There is an explicitly specified finite fully characteristic
quotient $D^*$ of $F/F^{(5)}$, of order divisible only by two and three,
such that the full character factors as
\[
 \Delta:\widehat{\mathrm{GT}}_0
       \longrightarrow \operatorname{Out}(D^*)
       \longrightarrow \mathbb F_2,
                                                        \tag{C.51.1}
\]
where the second arrow is defined on the actual finite outer image.
The construction and a bound use only the following integers:
\[
 \begin{split}
 M&=2^{293}3^{193},\qquad B=M^{M^2},\\
 R&=2^{2B}B,\qquad d_0=R^{R^2},\qquad a=2d_0,\\
 N&=2^{a+2}3^{a+1}.
 \end{split}                                             \tag{C.51.2}
\]
Writing $N_D=\ker(F\to D)$ and letting $N_{\rm ab}(N)$ be the
kernel of the canonical map to $(\mathbb Z/N)^2$, take
\[
 D^*=F/(N_D\cap N_{\rm ab}(N)).
                                                        \tag{C.51.3}
\]
Then
\[
 |D^*|\le d_0N^2,\qquad (D^*)^{(5)}=1.
                                                        \tag{C.51.4}
\]
This is a full-multiplier finite outer detector. It does not compute
the actual finite outer image, an efficient all-multiplier evaluation
table, or the smallest detector. The sharp derived depth four of
C.48 and its compactness construction remain separate.
\end{statement}

\begin{proof}[Complete candidate proof]
\paragraph{The finite source and its numerical bound.}
Use exactly the $L,C,G,D$ of C.50. There $|L|=M$,
$n=|C|\le B$, and
\[
 |G|=2^{2n/\exp(L)}n\le 2^{2B}B=R.
                                                        \tag{C.51.5}
\]
Consequently $|D|\le |G|^{|G|^2}\le d_0$. All prime divisors
of $|D|$ are two or three, and $D$ is generated by the two
marked images of $x,y$. An automorphism is determined by these
two images, so
\[
 |\operatorname{Out}(D)|
       \le|\operatorname{Aut}(D)|\le |D|^2\le d_0^2.
                                                        \tag{C.51.6}
\]
This is a numerical order bound. We do not infer that the order of
an automorphism group divides $|D|^2$, or bound its prime valuations
by those of $|D|$.

\paragraph{The actual arithmetic outer action is unramified away from two and three.}
Let
\[
 \rho_D:G_{\mathbb Q}\longrightarrow\operatorname{Out}(D)
                                                        \tag{C.51.7}
\]
be the actual outer action obtained from the geometric fundamental
group of $\mathbb P^1-\{0,1,\infty\}$ and its quotient $D$.
The kernel defining $D$ is fully characteristic, so this quotient
and this action are defined without choosing a Galois-stable
particular homomorphism to the auxiliary group $G$.

Fix a rational prime $q\ge5$. Every finite group in the construction
of $D$ has order prime to $q$, so the quotient map to $D$ factors
through the single maximal prime-to-$q$ geometric fundamental
group. Take the strict henselization of $\mathbb Z_q$. The coordinate
change
\[
                         z=\frac{2t}{t+1}
                                                        \tag{C.51.8}
\]
sends the three marked branches to $0,1,2$, which are pairwise
distinct on its residue fibre. All their pairwise intersection
multiplicities are therefore zero.

Apply Hasson--Yelton, Theorem 1.2(a),(b), printed p.~2,
to the constant complex power series $0,1,2$.
Their topological family is the product of a fixed three-punctured
sphere with a punctured disk, so its outer monodromy is trivial.
The cited theorem identifies this outer action with the actual
prime-to-$q$ algebraic outer action over the strictly henselian
field. Thus the inertia group $I_q$ acts trivially in
$\operatorname{Out}(D)$.
The same distinct-reduction case is explicitly recorded in that
paper's Example 3.4, printed p.~10.
\cite[Theorem~1.2 and Example~3.4]{PBC_HassonYelton2020}
Changing from the rational base point used in that theorem to our
tangential base point changes only an inner representative.
The conclusion concerns the outer action and is therefore unchanged.
In particular, this is a statement about the combined finite
two-three group $D$, not an inference from its separate pro-two
and pro-three images.

Let $E/\mathbb Q$ be the finite Galois extension fixed by
$\ker\rho_D$. We have proved
\[
 E/\mathbb Q\text{ is unramified outside }\{2,3\},
 \qquad [E:\mathbb Q]\le d_0^2.
                                                        \tag{C.51.9}
\]

\paragraph{An explicit conductor bound for the abelian intersection.}
Put $E_{\rm ab}=E\cap\mathbb Q^{\rm ab}$.
By Kronecker--Weber,
$\operatorname{Gal}(\mathbb Q^{\rm ab}/\mathbb Q)$ is identified
with $\widehat{\mathbb Z}^{\times}$ by the full cyclotomic
character. We use the theorem in
\cite[I, Theorem~4.16, pp.~49--50]{PBC_MilneCFT}.
The induced quotient to $\operatorname{Gal}(E_{\rm ab}/\mathbb Q)$
is trivial on $\mathbb Z_q^\times$ for every $q\ge5$:
this factor is the cyclotomic image of inertia at $q$, whereas
$E_{\rm ab}$ is unramified there.
Indeed, all prime-to-$q$ roots of unity are unramified locally and
the $q$-power cyclotomic extension is totally ramified.

The order of $\operatorname{Gal}(E_{\rm ab}/\mathbb Q)$ divides
$[E:\mathbb Q]$, and for $p=2,3$ we have
\[
 v_p([E:\mathbb Q])\le \log_p(d_0^2)\le 2d_0=a.
                                                        \tag{C.51.10}
\]
Hence its pro-two image is killed by $2^a$ and its pro-three
image is killed by $3^a$.
The elementary local unit decompositions give
\[
 \begin{split}
  \mathbb Z_2^\times&=\{\pm1\}\times(1+4\mathbb Z_2),
       & (1+4\mathbb Z_2)^{2^a}&=1+2^{a+2}\mathbb Z_2,\\
  \mathbb Z_3^\times&=\{\pm1\}\times(1+3\mathbb Z_3),
       & (1+3\mathbb Z_3)^{3^a}&=1+3^{a+1}\mathbb Z_3.
 \end{split}                                             \tag{C.51.11}
\]
For example, the power identities follow from the isomorphisms
with the additive $p$-adic integers given by the logarithm on the
displayed principal unit groups. They can also be checked using
$v_2(5^{2^k}-1)=k+2$ and
$v_3(4^{3^k}-1)=k+1$.

It follows that every full cyclotomic unit congruent to one modulo
$N$ acts trivially on $E_{\rm ab}$. Equivalently,
\[
                         E_{\rm ab}\subset\mathbb Q(\zeta_N).
                                                        \tag{C.51.12}
\]
The number $N$ need not be the conductor. It is the explicit upper
bound in (C.51.2); no calculation of the field $E$ is required.

\paragraph{The enlargement is an actual fully characteristic quotient.}
Both kernels in (C.51.3) are fully characteristic in $F$.
For the abelian one, this follows because every endomorphism acts
linearly on $F^{\rm ab}=\widehat{\mathbb Z}^2$ and preserves
$N\widehat{\mathbb Z}^2$. Thus $D^*$ is finite and fully
characteristic, with actual surjections
\[
              D^*\twoheadrightarrow D,\qquad
              D^*\twoheadrightarrow(\mathbb Z/N)^2.
                                                        \tag{C.51.13}
\]
Its joint map embeds it into $D\times(\mathbb Z/N)^2$.
This proves the order bound and derived-length assertion in
(C.51.4), as well as the restriction to the primes two and three.
All actual full gentle and Galois outer actions descend to this
same marked quotient.

\paragraph{Full arithmetic correction with the finite outer coordinate fixed.}
Suppose $g\in\widehat{\mathrm{GT}}_0$ has trivial outer action
on $D^*$. Its action on the abelian quotient in (C.51.13) is
the scalar $\lambda(g)$; an inner automorphism has trivial
abelian action. Therefore
\[
                         \lambda(g)\equiv1\pmod N.
                                                        \tag{C.51.14}
\]
Write $\lambda=\lambda(g)$ as an element of
$\widehat{\mathbb Z}^{\times}$. By (C.51.12), the prescribed
automorphisms $1$ of $E$ and $\lambda$ of $\mathbb Q^{\rm ab}$
agree on their intersection. Hence there is an actual
$\sigma\in G_{\mathbb Q}$ such that
\[
                   \rho_D(\sigma)=1,\qquad
                   \lambda(\sigma)=\lambda
                   \quad\hbox{in }\widehat{\mathbb Z}^{\times}.
                                                        \tag{C.51.15}
\]
Here the multiplier equality is exact at every prime, not just
modulo $N$. One can verify existence without any finite-level
lifting assumption: for each finite cyclotomic field $K$, the
restrictions $1$ on $E$ and $\lambda$ on $K$ agree on $E\cap K$.
They therefore define an automorphism of the finite compositum
$EK$ and extend to $\overline{\mathbb Q}$. The resulting nonempty
closed subsets of $G_{\mathbb Q}$ have the finite intersection
property as $K$ increases. Compactness gives (C.51.15)
simultaneously for all $K$.

View $\sigma$ in the specified arithmetic subgroup of the full
gentle group. The element
\[
                           h=g\sigma^{-1}
                                                        \tag{C.51.16}
\]
has multiplier one. Moreover $g$ has trivial outer action on
$D$, by the first quotient in (C.51.13), and so does $\sigma$.
Thus $h$ has trivial outer action on $D$.
The exact multiplier-one detector C.50 gives $\Delta(h)=0$.
The specified arithmetic subgroup lies in $\ker\Delta$, by
C.21. Since $\Delta$ is a character,
\[
                    \Delta(g)=\Delta(h)+\Delta(\sigma)=0.
                                                        \tag{C.51.17}
\]
We have proved containment of the full outer-action kernel in
$\ker\Delta$. This is precisely the factorization (C.51.1).
All maps involved have finite target after passage to $D^*$, so
the resulting character of the actual finite image is continuous.

No numerical enumeration or unknown full-multiplier frame formula
enters this proof. C.50 retains its explicit evaluation algorithm
on the multiplier-one image. The present argument extends the
existence of a specified finite outer detector to the whole group
by an actual arithmetic correction, without claiming that the
finite arithmetic image or a full-image evaluation table has been
computed.
\end{proof}
